%============================================================
% Chapter 29: Bit-Level Optimisation and Performance
% STATUS: PLANNED  (generated by tools/outline.py; replace this file when the chapter is written)
%============================================================
\chapter{Bit-Level Optimisation and Performance}
\label{ch:optimization}

\begin{tcolorbox}[enhanced, colback=white, colframe=secondaryblue, boxrule=1pt, arc=2mm]
\textbf{\large\textcolor{primaryblue}{Make bit-heavy code fast: built-ins and target pragmas, branchless code, masks instead of modulo, and memory packing.}}
\end{tcolorbox}

\chapterfacts{1800}{3000}{4}{8}{30}

\begin{plannedbox}[Planned Contents --- this chapter is not written yet]
\textbf{Part A --- Lessons}
\begin{enumerate}[leftmargin=*, label=\textbf{A\arabic*.}, itemsep=1pt]
    \item Compiler built-ins and \texttt{\#pragma GCC target("popcnt,bmi2")}
    \item Branchless min, max, abs and select with masks
    \item Power-of-two sizes: masking instead of modulo, ring buffers
    \item Packing data into bits to fit memory and cache
    \item Measuring: when a bit trick is worth it
\end{enumerate}
\textbf{Part B --- Patterns} (with candidate problems, checked when the chapter is written)
\begin{enumerate}[leftmargin=*, label=\textbf{P\arabic*.}, itemsep=1pt]
    \item \textbf{Built-ins}
    \item \textbf{Branchless code}
    \item \textbf{Packing}
    \item \textbf{Measuring}
\end{enumerate}
\end{plannedbox}
